\centerline{\bf \Huge Сравнения по модулю I}

\begin{flushright}
        \scriptsize{Работа тренера? Я не хочу терять футболиста в себе,\\ а это необходимо сделать, чтобы стать успешным тренером\\Деннис Бергкамп}
\end{flushright}

\textbf{\Large Теория}

\textbf{Определение.} Целые числа $a$ и $b$ называются \textit{сравнимыми по модулю} $m$, если они дают одинаковые остатки при делении на $m$. Запись: $a \equiv b\ (mod \ m)$.

\textbf{Важная теорема.} Для любых $a, b, \in \mathbb{Z} \ a \equiv b\ (mod \ m)$ тогда и только тогда, когда $a-b$ делится на $m$.

\textbf{Теорема.} Для любых $a, b, c$,

(a) Пусть $a \equiv b\ (mod \ m), c \equiv d\ (mod \ m)$, тогда $a \pm c \equiv b \pm d \ (mod \ m)$.

(b) Пусть $a \equiv b\ (mod \ m), c \equiv d\ (mod \ m)$, тогда $a c \equiv b d \ (mod \ m)$.

(c) Пусть $a \equiv b\ (mod \ m), k \in \mathbb{Z}$, тогда $k a \equiv k b\ (mod \ m)$.

(d) Пусть $a \equiv b\ (mod \ m), n \in \mathbb{N}$, тогда $a^n \equiv b^n\ (mod \ m)$.

(e) Пусть $k a \equiv k b\ (mod \ m), k \in \mathbb{Z}$ и $(k, m)=1$. Тогда $a \equiv b\ (mod \ m)$.

\textbf{\Large Практика}

1. Докажите, что число $1000 \cdot 1001 \cdot 1002 \cdot 1003-24$ делится на: (а) 999; (b) 1004.

2. Докажите без вычислений, что $1 \cdot 2 \cdot 3 \cdot \ldots \cdot 9+10 \cdot 11 \cdot \ldots \cdot 17 \cdot 18$ делится на 19.

3. Докажите, что число $9^{2023}+7^{2024}$ делится на 10 .

4. Докажите, что $1^{2025}+2^{2025}+\ldots+30^{2025}$ делится на 31 .

5. Докажите, что $2 \cdot 4 \cdot 6 \cdot 8 \cdot \ldots \cdot 2022 \cdot 2024-1 \cdot 3 \cdot 5 \cdot 7 \cdot \ldots \cdot 2021 \cdot 2023$ делится на 2025.

6. Найдите остаток от деления $2222^{5555}+5555^{2222}$ на 7.